\documentclass[11pt]{letter}
\usepackage[margin=1in]{geometry}
\usepackage{hyperref}
\hypersetup{colorlinks=true,linkcolor=black,urlcolor=blue}

\address{Andrew H.\ Bond\\
Department of Computer Engineering\\
San Jos\'{e} State University\\
\texttt{andrew.bond@sjsu.edu}\\
ORCID: 0009-0003-2599-6158}

\signature{Andrew H.\ Bond\\
Senior Lecturer, Department of Computer Engineering\\
San Jos\'{e} State University}

\begin{document}

\begin{letter}{Editor-in-Chief\\
\emph{Journal of Science and Law}\\
Center for Science and Law}

\opening{Dear Editor,}

I am pleased to submit the manuscript ``The Legal Bond Index: A
Geometric, Case-Level Test of Algorithmic Fairness Applied to COMPAS''
for consideration as an Original Research article in the
\emph{Journal of Science and Law}.

The paper addresses a problem that sits squarely within JSciLaw's
mission: a scientific result that bears directly on legal policy.
The 2016--2017 controversy over the COMPAS pretrial-risk algorithm
exposed a structural ambiguity in the concept of algorithmic
fairness, and the subsequent technical literature
(Chouldechova 2017; Kleinberg, Mullainathan and Raghavan 2017)
proved that the two competing fairness criteria each side invoked
are mathematically incompatible when group base rates differ. The
decade since has produced a sprawl of competing fairness metrics,
none of them dispositive, and a regulatory landscape in which courts
and oversight bodies are asked to evaluate algorithmic systems
without a single quantitative measure they can rely on.

This paper introduces such a measure: the \emph{Legal Bond Index}
(LBI), a matched-pair case-level test of whether an algorithm
treats individually-similar defendants similarly. LBI is rooted in
the Aristotelian like-cases-alike criterion, operationalized via the
Mahalanobis distance on legitimate predictors and a $k$-nearest-neighbour
matched-pair construction. LBI sits on a structurally distinct axis
from the calibration--equalized-odds debate: it does not constrain
population-level statistics, so the Chouldechova--Kleinberg
impossibility does not apply to it.

On the ProPublica COMPAS dataset ($n = 5{,}278$ defendants), the
empirical result is sharp. The standard fairness metrics that drove
the original controversy are reproduced exactly (FPR disparity
$+0.20$; calibration gap at high decile $+0.05$). The Legal Bond
Index returns $\mathrm{LBI}(\text{race}) = 1.050$ at neighbourhood
scale $k = 20$, with bootstrap 95\,\% CI $[1.040, 1.058]$ and
permutation $p < 0.005$. The 5\,\% case-level excess sensitivity to
race is small in absolute magnitude, robustly above the permutation
null centred at $1.00$ (the observed value sits 4 standard deviations
above the null), and replicates across six neighbourhood scales
spanning two orders of magnitude in $k$. It is a finding the
population-level metrics do not, and cannot, deliver.

The contributions are three:

\begin{enumerate}
\item A fairness criterion (LBI) that is case-level rather than
population-level, that returns a single scalar with a confidence
interval, and that requires only the algorithm's inputs and outputs
(not its parameters or training data)---a property essential for
black-box regulatory audits of proprietary systems.

\item An empirical demonstration on the canonical COMPAS dataset
that LBI produces a statistically robust, interpretable result that
neither calibration nor equalized odds alone can produce.

\item A reference implementation released as an open-source Jupyter
notebook, reproducible end-to-end from the public ProPublica file,
with explicit feature engineering, bootstrap CIs, and a
permutation-test certification of significance.
\end{enumerate}

The paper's relevance to \emph{JSciLaw}'s readership is direct.
Algorithmic risk assessment is in use in pretrial detention,
sentencing, parole, child-welfare, and lending decisions across
U.S.\ jurisdictions. The legal community needs a quantitative audit
tool that produces a single number, can be applied to black-box
algorithms, and is grounded in a defensible interpretation of equal
treatment. LBI fills that gap. The paper also bridges scientific
method and legal doctrine: the metric is geometric (Mahalanobis
distance, $k$-nearest neighbours, matched-pair design) and the
interpretation is doctrinal (Aristotle's like-cases-alike;
individual fairness in Dwork et al.\ 2012's sense).

The manuscript is self-contained. The broader theoretical framework
(``Judicial Complex,'' gauge-theoretic structure, path-homology
criteria for constitutional review) that motivates the LBI is
mentioned in Section 4 as one paragraph; readers who care about the
mathematical foundations are pointed to the companion theoretical
work; readers who care only about the audit tool need no part of
that machinery.

Data, code, and figures are reproducible. The COMPAS file is public.
The Jupyter notebook (\verb|lbi_compas_analysis.ipynb|) reproduces
all tables and figures from the public CSV in approximately five
minutes on a laptop. The reference implementation is bundled with
the submission and will be released on GitHub on acceptance.

The manuscript has not been published previously and is not under
consideration elsewhere. There are no competing interests and no
external funding.

I would be glad to suggest reviewers with expertise in algorithmic
fairness, computational law, and criminal-justice risk assessment if
that would help expedite review.

\closing{Sincerely,}

\end{letter}

\end{document}
